% Even and odd alignments of one ordered tetrahedral ligand frame.
% Input-only TikZ file; requires concept_style.tex (Okabe--Ito palette, \iosub)
% and the libraries arrows.meta, calc, positioning, backgrounds.
% Ligand colors follow the shared permutation gallery (orbit_examples_six.tex):
% ligand 1 blue, 2 orange, 3 green, 4 purple (Paul Tol qualitative, CVD-safe).
% Numbers inside the balls name ligands; small gray digits name frame slots.
\providecolor{pballA}{RGB}{68,119,170}
\providecolor{pballB}{RGB}{238,119,51}
\providecolor{pballC}{RGB}{34,136,51}
\providecolor{pballD}{RGB}{170,51,119}

\begin{tikzpicture}[
  font=\small,
  bball/.style={circle, inner sep=0pt, line width=0.6pt,
                minimum size=6.4mm, font=\bfseries\footnotesize},
  sball/.style={circle, inner sep=0pt, line width=0.5pt,
                minimum size=4.6mm, font=\bfseries\scriptsize},
  lig1/.style={fill=pballA, draw=pballA!62!black, text=white},
  lig2/.style={fill=pballB, draw=pballB!62!black, text=black!88},
  lig3/.style={fill=pballC, draw=pballC!62!black, text=white},
  lig4/.style={fill=pballD, draw=pballD!62!black, text=white},
  slotn/.style={font=\scriptsize, text=black!55, inner sep=0.6pt},
  state/.style={font=\footnotesize, text=black!82},
  tag/.style={font=\scriptsize, text=black!50},
  evenarr/.style={-{Latex[length=2.2mm,width=1.6mm]}, line width=0.9pt,
                  draw=ioGreen!62!black},
  oddarr/.style={-{Latex[length=2.2mm,width=1.6mm]}, line width=0.9pt,
                 draw=ioBlue!72!black},
  veq/.style={rounded corners=2pt, draw=ioGreen!60!black, fill=ioGreen!14,
              font=\footnotesize\bfseries, text=ioGreen!45!black, inner sep=2.4pt},
  vopp/.style={rounded corners=2pt, draw=ioBlue!55!black, fill=ioBlue!10,
               font=\footnotesize\bfseries, text=ioBlue!45!black, inner sep=2.4pt},
  chip/.style={rounded corners=2pt, font=\footnotesize, inner sep=2.6pt}
]

% ---- specular highlights that turn a flat disc into a 3D sphere -----------
\newcommand{\bpHlB}[2]{%
  \fill[white,opacity=0.18] ({#1-0.096},{#2+0.096}) circle (0.096);}
\newcommand{\bpHlS}[2]{%
  \fill[white,opacity=0.18] ({#1-0.069},{#2+0.069}) circle (0.069);}

% ---- tetrahedral center: two plain bonds, one wedge (front), one hash -----
% #1,#2 center; #3..#6 = ligand style / ligand label for slots 1..4.
\newcommand{\bpTetra}[6]{\begin{scope}[shift={(#1,#2)}]
  \draw[iobond] (0,0) -- (-0.80,0.62);
  \draw[iobond] (0,0) -- (0.80,0.62);
  % slot 3: bold wedge toward the viewer
  \fill[black!82] (0,0) -- (0.826,-0.584) -- (0.694,-0.736) -- cycle;
  % slot 4: hashed wedge away from the viewer
  \foreach \t/\h in {0.16/0.024,0.30/0.036,0.44/0.048,0.58/0.060,0.70/0.070}
    \draw[line width=0.7pt, black!82]
      ({-0.76*\t+0.656*\h},{-0.66*\t-0.755*\h}) --
      ({-0.76*\t-0.656*\h},{-0.66*\t+0.755*\h});
  \fill[black!78] (0,0) circle (0.055);
  \foreach \px/\py/\sty/\lab/\slt/\lx/\ly in {%
      -0.80/0.62/#3/1/-1.18/0.84, 0.80/0.62/#4/2/1.18/0.84,
      0.76/-0.66/#5/3/1.14/-0.88, -0.76/-0.66/#6/4/-1.14/-0.88}{
    \node[bball,\sty] at (\px,\py) {\lab};
    \bpHlB{\px}{\py}
    \node[slotn] at (\lx,\ly) {\slt};}
\end{scope}}

% ---- the same state as an ordered frame: values in slot order -------------
\newcommand{\bpStrip}[6]{\begin{scope}[shift={(#1,#2)}]
  \draw[iocard] (-1.16,-0.46) rectangle (1.16,0.30);
  \foreach \px/\sty/\lab/\slt in {%
      -0.87/#3/1, -0.29/#4/2, 0.29/#5/3, 0.87/#6/4}{
    \node[sball,\sty] at (\px,0) {\lab};
    \bpHlS{\px}{0}
    \node[slotn] at (\px,-0.31) {\slt};}
\end{scope}}

% ============================ Panel A: even ================================
\begin{scope}[shift={(0,0)}]
  \iosub[{$\epsilon\in\Gamma_{\mathrm{tet}}$; $12$ of the $24$ alignments}]%
    {-0.42,0.30}{A}{Even action \;\;$\epsilon\in A_4=\Gamma_{\mathrm{tet}}$}
  \node[state] at (1.60,-0.86) {$u=(1,2,3,4)$};
  \bpTetra{1.60}{-2.05}{lig1/1}{lig2/2}{lig3/3}{lig4/4}
  \bpStrip{1.60}{-3.55}{lig1/1}{lig2/2}{lig3/3}{lig4/4}
  \draw[evenarr] (2.98,-2.05) -- (4.02,-2.05);
  \node[state] at (3.50,-1.80) {$\epsilon=(1\;2\;3)$};
  \node[tag]   at (3.50,-2.32) {$\operatorname{sgn}\epsilon=+1$};
  \node[state] at (5.40,-0.86) {$\epsilon\!\cdot\!u=(2,3,1,4)$};
  \bpTetra{5.40}{-2.05}{lig2/2}{lig3/3}{lig1/1}{lig4/4}
  \bpStrip{5.40}{-3.55}{lig2/2}{lig3/3}{lig1/1}{lig4/4}
  \node[veq] at (3.50,-4.45)
    {SAME ORBIT \;{\normalfont\footnotesize
     $[\epsilon\!\cdot\!u]_{A_4}=[u]_{A_4}$}};
\end{scope}

% ============================ Panel B: odd =================================
\begin{scope}[shift={(8.60,0)}]
  \iosub[{$\rho\in S_4\setminus\Gamma_{\mathrm{tet}}$; $12$ of the $24$}]%
    {-0.42,0.30}{B}{Odd action \;\;$\rho\notin A_4$}
  \node[state] at (1.60,-0.86) {$u=(1,2,3,4)$};
  \bpTetra{1.60}{-2.05}{lig1/1}{lig2/2}{lig3/3}{lig4/4}
  \bpStrip{1.60}{-3.55}{lig1/1}{lig2/2}{lig3/3}{lig4/4}
  \draw[oddarr] (2.98,-2.05) -- (4.02,-2.05);
  \node[state] at (3.50,-1.80) {$\rho=(1\;2)$};
  \node[tag]   at (3.50,-2.32) {$\operatorname{sgn}\rho=-1$};
  \node[state] at (5.40,-0.86) {$\rho\!\cdot\!u=(2,1,3,4)$};
  \bpTetra{5.40}{-2.05}{lig2/2}{lig1/1}{lig3/3}{lig4/4}
  \bpStrip{5.40}{-3.55}{lig2/2}{lig1/1}{lig3/3}{lig4/4}
  \node[vopp] at (3.50,-4.45)
    {OPPOSITE \;{\normalfont\footnotesize
     $[\rho\!\cdot\!u]_{A_4}=\operatorname{opp}[u]_{A_4}$}};
\end{scope}

% ---- coset ledger shared by both panels -----------------------------------
\node[chip, draw=ioGreen!55!black, fill=ioGreen!12, text=ioGreen!42!black]
  at (1.75,-5.52) {$[u]_{A_4}$: $12$ even frames};
\node[font=\footnotesize, text=black!70, align=center] at (7.80,-5.52)
  {$S_4=A_4\,\sqcup\,\rho A_4$, hence $24=12+12$:\\
   $[S_4:A_4]=2$ configurations, exchanged by every odd alignment};
\node[chip, draw=ioBlue!55!black, fill=ioBlue!10, text=ioBlue!42!black]
  at (13.85,-5.52) {$[\rho\!\cdot\!u]_{A_4}$: $12$ odd frames};

% ---- panel and ledger backgrounds -----------------------------------------
\begin{pgfonlayer}{background}
  \foreach \x in {0,8.60}
    \draw[iopanel] ($(\x,0)+(-0.55,-4.84)$) rectangle ($(\x,0)+(7.55,0.55)$);
  \draw[iopanel] (-0.55,-6.02) rectangle (16.15,-5.02);
\end{pgfonlayer}
\end{tikzpicture}
